% Generated draft: requires mathematical and source review.
\documentclass[11pt]{article}
\usepackage[margin=25mm]{geometry}
\usepackage{fontspec}
\setmainfont{Latin Modern Roman}
\usepackage{amsmath,amssymb,mathtools,unicode-math}
\setmathfont{Latin Modern Math}
\usepackage{enumitem,needspace,xurl,hyperref}
\setlength{\parindent}{0pt}
\setlength{\parskip}{5pt}
\setlength{\emergencystretch}{3em}
\newcommand{\N}{\mathbb{N}}
\newcommand{\Z}{\mathbb{Z}}
\newcommand{\Q}{\mathbb{Q}}
\newcommand{\R}{\mathbb{R}}
\newcommand{\C}{\mathbb{C}}
\newcommand{\F}{\mathbb{F}}
\newcommand{\A}{\mathbb{A}}
\newcommand{\PP}{\mathbb{P}}
\newcommand{\E}{\mathbb{E}}
\newcommand{\eps}{\varepsilon}
\newcommand{\dd}{\,\mathrm{d}}
\begin{document}
\section*{UnsolvedMath ID: 10400161}
\textbf{Source problem code:} catalogue entry
\subsection{Definitions and mathematical statement}
The source's \(q/l\)-solvable and Casson-type TQFT terminology is not defined in the excerpt. Retaining these as source-specific predicates, Kerler's program has six parts:
\begin{enumerate}
\item Lift the indicated \(1/1\)-solvable Casson-type TQFTs over \(\mathbb F_p\) to a universal integral \(1/1\)-solvable Casson-type theory over \(\mathbb Z\).
\item Describe its resulting invariant \(\Xi_{\mathbb Z}(M)\) for three-manifolds with \(b_1(M)\ge1\), where \(b_1=\dim_{\mathbb Q}H_1(M;\mathbb Q)\).
\item Develop perturbation theory for general \(q/l\)-solvable TQFTs.
\item Relate these theories to the specified standard resolutions of \(\Gamma_g\).
\item Relate them to traditional finite-type theory of closed three-manifolds.
\item Describe Reshetikhin--Turaev theories within the same pattern.
\end{enumerate}
The parameter ranges, prime restrictions, universal property and compatibility of reduction modulo \(p\) are not supplied. \(\Gamma_g\) is the intended mapping-class-group notation, but its boundary conventions and the resolutions meant are absent. Neither \(\Xi_{\mathbb Z}\) nor Casson type is replaced by the ordinary Casson invariant, whose usual domain differs from the stated \(b_1\ge1\) request.

The passage supplies categorical context. Its \(\operatorname{Cob}_\bullet\) has compact oriented one-boundary-component surfaces as objects and suitable three-cobordisms, with a braided tensor structure. The algebraic category \(\operatorname{Alg}\) is generated on objects by \(A\), so objects are \(A^{\otimes g}\), with unit \(A^{\otimes0}\). Morphisms are generated under composition and tensor product by the ribbon-Hopf-algebra structure, including
\[
 m:A\otimes A\to A,\quad\Delta:A\to A\otimes A,\quad
 e:\mathbf1\to A,\quad\varepsilon:A\to\mathbf1,
\]
with integrals and a nondegenerate pairing. The full relations are not given. The source reports a surjective functor \(\operatorname{Alg}\to\operatorname{Cob}_\bullet\), an isomorphism on its genus-one restriction, in reference [215]. It reports preparatory integral Casson formulas in [212] and calls item (2) immediate from preceding discussion, which is not supplied. None of these contextual claims is independently proved here.

The summary calls the program not obviously resolved. This remains a qualified source assessment, not a verified current solution or a completed definition of the six tasks.
\subsection{Short English statement}
Develop the integral lift, invariant, perturbation theory and mapping-class/finite-type/Reshetikhin--Turaev comparisons in Kerler's solvable-TQFT program. All six tasks and the Hopf/cobordism context are retained, but the solvability and universality definitions and referenced constructions are missing.
\subsection{Sources}
\begin{itemize}
\item Local registry: \url{data/unsolvedmath/problems.json}, permanent UnsolvedMath ID \(10400161\), source code catalogue entry; complete supplied statement, background, and research summary, updated 2026-08-06.
\item Ohtsuki, \emph{Problems on invariants of knots and 3-manifolds} (2002), Problem 8.15, attributed to T. Kerler, PDF page 130: \url{https://msp.org/gtm/2002/04/gtm-2002-04-024s.pdf}. Supplied PDF-font-structure extraction accessed 2026-07-29.
\item Internal reference [212] is cited for integral Casson formulas; [215] for the algebraic-to-cobordism functor and genus-one result. Their bibliographic entries are not included.
\end{itemize}
\end{document}
